\scanpage{200}

Si on veut étendre $\varphi_n^0$ en \deleted{$\varphi_n$} : hom de $S$-algèbre $\varphi_n:\mathcal A_n\unsure{-|\to}\mathcal A_n$, il faut se donner les images des $dU_0,-,dU_{n-1}$, et vérifier seulement que ce sont des éléments de \uline{carré nul}, \deleted{\unsure{de degré impair}} \inserted{qui anticomm$^{\mathrm t}$ \deleted{\unsure{deux à deux}} entre eux et commutent} aux $\xi_{s,i}$

\[
\varphi_n(dU_r)=d\varphi_n(U_r)=d\sigma_r^*(\xi_1)=\sigma_r^*(d\xi_1)
\]

\[
\begin{aligned}
 d\xi_1&=\sum_{j\geq1}a_{1j}X_0^{(j-1)}dX_0
 &\biggl(&=d\Bigl(\sum_{j\geq0}a_{1j}U_r^{(j)}\Bigr)\\
 \deleted{d\xi_1}&\deleted{=\sum_{j\geq1}a_{1j}X_0^{(j-1)}dX_0}
 &&=\Bigl(\sum_{j\geq1}a_{1j}U_r^{(j-1)}\Bigr)dU_r\biggr)
\end{aligned}
\]

Ces éléments sont bien de carré nul, anticommutent entre eux, mais commutent-ils aux $\xi_{s,i}=\sigma_s^*(\deleted{\unsure{\xi}}X_0^{(i)})=\sum_{j\geq0}a_{ij}U_s^{(j)}+(\sum_{j\geq0}b_{ij}U_s^{(j)})dU_s$

\deleted{\unsure{En résumé, écrivons|En résumé, écrivant} $\mathcal A_n=S\{U_0,-,U_{n-1}\}[dU_0,-,dU_{n-1}]$ et pour l'élém\textsuperscript{t} général $=F+\sum_{0\leq r\leq n-1}G_r\,dU_r$.}

% slanted
\margin{et $\sigma_r^*(d\xi_1)$ commute bien aux $\xi_{si}=\sigma_s^*(X_0^{(i)})$ pour $s=r$, mais pour $s\ne r$ c'est vrai ssi $2\bigl(\sum_{j\geq1}a_{1j}U_r^{(j-1)}\bigr)\bigl(\sum_{j\geq0}b_{ij}U_s^{(j)}\bigr)=0$, i.e. $\boxed{\substack{2a_{1,j}b_{i\deleted{j}k}=0\\\forall j\geq1,\ k\geq0}}$ Or $\unsure{\alpha}a_{11}=1+\sum_{i\geq\deleted{\unsure{0}}1}\alpha_i T^{(i)}$ est \unsure{non|un} div de \unsure{Zéro|Zéros}, donc il faut $2b_{ik}=0\quad(\forall i\geq1,k\geq0)$ donc $2B_i=0$.}

On utilise, pour prouver que $\varphi_n$ (qui à priori est un hom de $S$-algèbres) commute à la différentielle, on utilise le

\uline{Lemme} $S$ anneau commutatif, $U_0-U_m$ des variables, on considère la $S$-algèbre diff $\mathcal A\deleted{\unsure{\text{anneau}}}=S\{U_0,-,U_m\}[dU_0,-,dU_m]$ (\uline{NB} on néglige la structure des p.d, \unsure{au moins|du moins} $dU_r^{(i)}=U_r^{(i-1)}dU_r$ --). Soit $B$ une $S$-alg. diff. \unsure{alternée}, (grad. mod 2, différentielle de degré $+1$), On s'intéresse aux hom de \inserted{$S$-}algèbres diff $\mathcal A\xrightarrow{\varphi} B$. Pour les \unsure{dét.|tels|tel}, on a

% archive sheet number: 199
